\section{动量与动量守恒}

\subsection{动量与冲量}
我们将质点的速度的$\vb*{v}$与其质量$m$的乘积，定义为该质点的\uwave{动量}，记作$\vb*{P}$
\begin{BoxDefinition}[动量]
    \vb*{P}=m\vb*{v}
\end{BoxDefinition}
动量是与参照系有关的量，因为速度会随参照系而变。

\Refx{定义}{动量}是描述质点运动状态强弱的矢量，动量的方向与速度相同，动量的大小同时与速度和质量相关，因为生活经验告诉我们，玩具小车在速度较低时是无需担心的，但是在相同速度下，集装箱卡车则是非常危险的，由此可见质量和运动的强弱是密切相关的。

\Refx{定义}{动量}先前已经在\Refx{公式}{牛顿第二定律的动量形式}中出现过，该公式告诉我们，所谓改变物体运动状态的力，实际上就是动量的变化率。

如果有一个质点受到恒力$\vb*{F}$的作用，且这种作用持续了$\delt{t}$的时间，那么我们将力$\vb*{F}$与作用时间$\delt{t}$的乘积，称为力$\vb*{F}$产生的\uwave{冲量}，记作$\vb*{I}$，这可以表示为
\begin{BoxDefinition}[恒力的冲量]
    \vb*{I}=\vb*{F}\delt{t}
\end{BoxDefinition}

如果力$\vb*{F}$是随时间变化的变力，那么冲量可以用定积分表示
\begin{BoxDefinition}[变力的冲量]
    \vb*{I}=\Int[t_a][t_b]{\vb*{F}\dif{t}}
\end{BoxDefinition}
冲量是与参照系无关的绝对量，因为力是绝对的。

根据\Refx{定义}{变力的冲量}可知，\textbf{冲量的物理意义是力在时间上的累积}。

我们知道，力是动量的变化率，力对时间的积累即动量的变化率对时间的积累。

从而在\Refx{定义}{变力的冲量}中代入\Refx{公式}{牛顿第二定律的动量形式}
\begin{equation*}
    \vb*{I}=\Int[t_a][t_b]{\vb*{F}\dif{t}}=\Int[t_a][t_b]{\Fdif{\vb*{P}}{t}\dif{t}}=\Int[\vb*{P}_a][\vb*{P}_b]{\dif{\vb*{P}}}=\vb*{P}_b-\vb*{P}_a
\end{equation*}
这就得到了\uwave{动量定理}
\begin{BoxPrinciple}[动量定理]
    合力产生的冲量，等于质点的动量增量。
    \begin{BoxEq}
        \vb*{I}=\vb*{P}_b-\vb*{P}_a=\delt{\vb*{P}}
    \end{BoxEq}
\end{BoxPrinciple}
\Refx{定律}{动量定理}是\Refx{公式}{牛顿第二定律的动量形式}的应用结果，因为力是动量的变化率，因此冲量，即力在时间上的积累，必然是动量的变化量。

动量不是绝对量，其会随着参照系的变化而变化，但是对于任意两个惯性系，由于参照系间的相对速度为常数，相同时刻的动量间仅相差一个常量。即在两个惯性系中的同一时间段，动量的变化量是恒定的，冲量又是绝对的，因此\textbf{动量定理普遍适用于一切惯性系}。

动量和冲量在国际单位制中的单位是\si{N\cdot s=kg\cdot m\cdot s^{-1}}，量纲是\si{[LMT^{-1}]}。\vspace{10pt}

\subsection{质点系的动量定理}
设两个质点的质量分别为$m_1,m_2$，在$t_a$至$t_b$的时刻间
\begin{enumerate}
    \item 两者分别受到外力$\vb*{F}_1,\vb*{F}_2$和内力$\vb*{f}_1,\vb*{f}_2$的作用。
    \item 两者受到的外力的冲量为$\vb*{I}_{1,E},\vb*{I}_{2,E}$，两者受到的内力的冲量为$\vb*{I}_{1,\vb*{I}},\vb*{I}_{2,\vb*{I}}$。
\end{enumerate}

% 分别受到外力$\vb*{F}_1,\vb*{F}_2$和内力$\vb*{f}_1,\vb*{f}_2$的作用，两者受到的外力的冲量为$\vb*{I}_{1,E},\vb*{I}_{2,E}$，两者受到的内力的冲量为$\vb*{I}_{1,\vb*{I}},\vb*{I}_{2,\vb*{I}}$。

对质点$m_1$运用\Refx{定律}{动量定理}，设$t_a,t_b$时其动量分别为$\vb*{P}_{1,a},\vb*{P}_{1,b}$
\begin{Equation}[质点系动量定理1]
    \vb*{I}_1=\vb*{I}_{1,E}+\vb*{I}_{1,I}=\Int[t_a][t_b]{(\vb*{F}_1+\vb*{f}_1)\dif{t}}=\vb*{P}_{1,b}-\vb*{P}_{1,a}
\end{Equation}
对质点$m_2$运用\Refx{定律}{动量定理}，设$t_a,t_b$时其动量分别为$\vb*{P}_{2,a},\vb*{P}_{2,b}$
\begin{Equation}[质点系动量定理2]
    \vb*{I}_2=\vb*{I}_{2,E}+\vb*{I}_{2,I}=\Int[t_a][t_b]{(\vb*{F}_2+\vb*{f}_2)\dif{t}}=\vb*{P}_{2,b}-\vb*{P}_{2,a}
\end{Equation}
上两式成立，是因为质点所受合力的冲量，等同于质点所受内力和外力的冲量和。
\begin{itemize}
    \item 外力对质点系的冲量，等同于外力对各质点的冲量和，即$\vb*{I}_E=\vb*{I}_{1,E}+\vb*{I}_{2,E}$。
    \item 内力对质点系的冲量，等同于内力对各质点的冲量和，即$\vb*{I}_I\hspace{2.2pt}=\vb*{I}_{1,I}\hspace{2.2pt}+\vb*{I}_{2,I}\hspace{2.2pt}$。
    \item 质点系的动量，等同于各质点的动量和，即$\vb*{P}_{a}=\vb*{P}_{1,a}+\vb*{P}_{2,a}$和$\vb*{P}_{b}=\vb*{P}_{1,b}+\vb*{P}_{2,b}$。
\end{itemize}
但是关注到，内力$\vb*{f}_1,\vb*{f}_2$是$m_1,m_2$间的一对相互作用力，根据\Refx{定律}{牛顿第三定律}
\begin{equation*}
    \vb*{I}_{I}=\vb*{I}_{1,I}+\vb*{I}_{2,I}=\Int[t_a][t_b]{(\vb*{f}_1+\vb*{f}_2)\dif{t}}=\vb*{0}
\end{equation*}
即质点系合内力的冲量$\vb*{I}_I$恒为零矢量，因而式\ref{式:质点系动量定理1}和式\ref{式:质点系动量定理2}相加得到
\begin{align*}
    \vb*{I}_E
    &=\qty(\vb*{P}_{1,b}-\vb*{P}_{1,a})+\qty(\vb*{P}_{2,b}-\vb*{P}_{2,a})\\
    &=\qty(\vb*{P}_{1,b}+\vb*{P}_{2,b})-\qty(\vb*{P}_{1,a}+\vb*{P}_{1,b})\\
    &=\vb*{P}_{b}-\vb*{P}_{a}
\end{align*}\vspace{-20pt}
\begin{BoxPrinciple}[质点系的动量定理]
    质点系的外力的冲量的矢量和，等于质点系的动量增量。
    \begin{BoxEq}
        \vb*{I}_E=\vb*{P}_{b}-\vb*{P}_{a}=\delt{\vb*{P}}
    \end{BoxEq}
\end{BoxPrinciple}
\Refx{定律}{质点系的动量定理}与先前的\Refx{定律}{质点系的动能定理}不同
\begin{itemize}
    \item 内力不能改变动量，因为动量是矢量的，因此内力导致的动量增量相互抵消。
    \item 内力可以改变动能，因为动能是标量的，因此内力导致的动能增量相互叠加。
\end{itemize}
\Refx{定律}{质点系的动量定理}只能在惯性系中使用，在非惯性系中不成立。

\subsection{动量守恒定律}
\Refx{定律}{质点系的动量定理}中，如果系统不受外力，关注到动量不变。

这就得到了以下的\uwave{动量守恒定律}
\begin{BoxPrinciple}[动量守恒定律]
    当质点系是孤立的，即所受的合外力为零，各质点仅受内力作用时，

    那么，各质点的动量可以相互转化，但是系统的动量总和保持不变，称为动量守恒。
\end{BoxPrinciple}
动量守恒定律要求质点系所受的合外力为零，但是当系统的内力远大于系统的外力时，例如在碰撞和爆炸的过程中，系统所受的外力可以近似为零，从而也适用动量守恒。

\section{碰撞}

\subsection{碰撞的概念}
如果两个或几个物体相互接近或发生接触时，在极为短暂的时间内，使物体的运动状态发生显著变化，那么将这个作用过程称为\uwave{碰撞}。通常将相互碰撞的物体作为一个系统考虑，由于碰撞时物体间的内力远大于外力，故碰撞系统应遵从\Refx{定律}{动量守恒定律}。

我们以两球碰撞为例进行讨论，若两球碰撞前后其速度均在两球中心连线上，那么，我们将这种碰撞称为\uwave{对心碰撞}或\uwave{正碰撞}，否则称之为\uwave{斜碰撞}，在这里我们只讨论对心碰撞的情况。

我们知道，碰撞过程的物体相互作用是极为复杂的，通过受力分析碰撞过程的结果是不可行的，设小球的质量$m_1,m_2$，已知碰撞前的速度为$v_{1a},v_{2a}$，试球碰撞后的速度为$v_{1b},v_{2b}$。

根据\Refx{定律}{动量守恒定律}，满足
\begin{Equation}[碰撞的动量守恒]
    m_1v_{1a}+m_2v_{2a}=m_1v_{1b}+m_2v_{2b}
\end{Equation}
但关注到，仅依靠式\ref{式:碰撞的动量守恒}不足以解出碰撞后的速度$v_{1b},v_{2b}$，因此在动量守恒之外还有方程。

牛顿通过实验总结了以下的\uwave{碰撞定律}（注意两个速度的相减方向相反）
\begin{itemize}
    \item 记碰撞前的速度差$(v_{1a}-v_{2a})$为两球的\uwave{接近速度}。
    \item 记碰撞后的速度差$(v_{2b}\hspace{0.7pt}-v_{1b}\hspace{0.7pt})$为两球的\uwave{分离速度}。
\end{itemize}
碰撞定律可以表述为，两球的分离速度与接近速度成正比，比值由材料性质决定，比值称为两球的\uwave{恢复系数}，记作$e$，其中$e\in[0,1]$，该定律可以写成以下的公式
\begin{BoxEquation}[碰撞定律]
    e=\frac{v_{2b}-v_{1b}}{v_{1a}-v_{2a}}
\end{BoxEquation}
联立\Refx{公式}{碰撞定律}和式\ref{式:碰撞的动量守恒}，得到关于$v_{2b},v_{1b}$的方程组
\begin{equation*}
    \begin{cases}
        v_{2b}-v_{1b}=ev_{1a}-ev_{2a}\\
        m_2v_{2b}+m_1v_{1b}=m_1v_{1a}+m_2v_{2a}
    \end{cases}
\end{equation*}
列出相应的行列式
\begin{align*}
    &D=
    \begin{vmatrix}
        1&-1\\
        m_2&m_1
    \end{vmatrix}=m_1+m_2\\[4mm]
    &D_{v_{1b}}=
    \begin{vmatrix}
        1&ev_{1a}-ev_{2a}\\
        m_2&m_1v_{1a}+m_2v_{2a}
    \end{vmatrix}=m_1v_{1a}+m_2v_{2a}-m_2e(v_{1a}-v_{2a})\\[4mm]
    &D_{v_{2b}}=
    \begin{vmatrix}
        ev_{1a}-ev_{2a}&-1\\
        m_1v_{1a}+m_2v_{2a}&m_1
    \end{vmatrix}=m_1v_{1a}+m_2v_{2a}+m_1e(v_{1a}-v_{2a})
\end{align*}
根据线性方程组理论
\begin{align*}
    v_{1b}=&\frac{D_{v1b}}{D}=
    \frac{m_1v_{1a}+m_2v_{2a}-m_2e(v_{1a}-v_{2a})}{m_1+m_2}\\[3mm]
    \xlongequal{\text{加减$m_2v_{1a}$}}&
    \frac{(m_1+m_2)v_{1a}-m_2(v_{1a}-v_{2a})-m_2e(v_{1a}-v_{2a})}{m_1+m_2}
    =v_{1a}-\frac{m_2(1+e)(v_{1a}-v_{2a})}{m_1+m_2}\\[8mm]
    v_{2b}=&\frac{D_{v1b}}{D}=
    \frac{m_1v_{1a}+m_2v_{2a}+m_1e(v_{1a}-v_{2a})}{m_1+m_2}\\[3mm]
    \xlongequal{\text{加减$m_1v_{2a}$}}&
    \frac{(m_1+m_2)v_{2a}+m_1(v_{1a}-v_{2a})+m_1e(v_{1a}-v_{2a})}{m_1+m_2}
    =v_{2a}+\frac{m_1(1+e)(v_{1a}-v_{2a})}{m_1+m_2}
\end{align*}
这就得到了对心碰撞在一般情况的$v_{1b},v_{2b}$的计算公式，整理如下。
\begin{BoxEquation}[一般对心碰撞的计算公式]
    \begin{cases}
        \mathlarger{v_{1b}=
        \frac{m_1v_{1a}+m_2v_{2a}-m_2e(v_{1a}-v_{2a})}{m_1+m_2}=v_{1a}-\frac{m_2(1+e)(v_{1a}-v_{2a})}{m_1+m_2}}\\[6mm]
        \mathlarger{v_{2b}=
        \frac{m_1v_{1a}+m_2v_{2a}+m_1e(v_{1a}-v_{2a})}{m_1+m_2}=v_{2a}+\frac{m_1(1+e)(v_{1a}-v_{2a})}{m_1+m_2}}
    \end{cases}
\end{BoxEquation}

\subsection{完全弹性碰撞}
当恢复系数$e=1$时，此时分离速度与接近速度相同，称为\uwave{完全弹性碰撞}。

在\Refx{公式}{一般对心碰撞的计算公式}中的左式，令$e=1$得
\begin{align*}
    v_{1b}&=\frac{m_1v_{1a}+m_2v_{2a}-m_2(v_{1a}-v_{2a})}{m_1+m_2}=\frac{(m_1-m_2)v_{1a}+2m_2v_{2a}}{m_1+m_2}\\[4mm]
    v_{2b}&=\frac{m_1v_{1a}+m_2v_{2a}+m_1(v_{1a}-v_{2a})}{m_1+m_2}=\frac{(m_2-m_1)v_{2a}+2m_1v_{1a}}{m_1+m_2}  
\end{align*}
这就得到了完全弹性碰撞的$\vb*{v}_{1b},\vb*{v}_{2b}$的计算公式，整理如下。
\begin{BoxEquation}[完全弹性碰撞的计算公式]
    \begin{cases}
        \mathlarger{v_{1b}=\frac{m_1v_{1a}+m_2v_{2a}-m_2(v_{1a}-v_{2a})}{m_1+m_2}=\frac{m_1-m_2}{m_1+m_2}v_{1a}+\frac{2m_2}{m_1+m_2}v_{2a}}\\[6mm]
        \mathlarger{v_{2b}=\frac{m_1v_{1a}+m_2v_{2a}+m_1(v_{1a}-v_{2a})}{m_1+m_2}=\frac{m_2-m_1}{m_1+m_2}v_{2a}+\frac{2m_1}{m_1+m_2}v_{1a}}
    \end{cases}
\end{BoxEquation}
下面，我们分析完全弹性碰撞中的几种特例
\begin{enumerate}
    \item 设两球的质量相等，即$m_2=m_1$时
    \begin{equation*}
        v_{1b}=v_{2a}\qquad v_{2b}=v_{1a}
    \end{equation*}
    即两球质量相同时，碰撞时将交换彼此交换速度，这个现象称为速度的\uwave{转移}。特别的，如果第二个小球原为静止$v_{2a}=0$，那么当第一个小球在以速度$v_{1a}$与其发生碰撞时，第一个小球就将停止运动，并将$v_{1a}$的速度传递给第二个小球，这就是\uwave{牛顿摆}。

    \item 设第二个物体质量极大且保持静止，即$m_2\gg m_1$且$v_{2a}=0$时
    \begin{equation*}
        v_{1b}=-v_{1a}\qquad v_{2b}=0
    \end{equation*}
    即静止的质量极大的物体在碰撞时仍基本保持静止，而质量较小的物体在发生碰撞后，速度反向，速度大小基本不变，这样的现象称为\uwave{反冲}，小球撞击墙壁就是反冲的例子。
\end{enumerate}
我们知道，对于一般碰撞的而言，尽管动量守恒，但动能未必守恒。

现在我们要证明，完全弹性碰撞事实上是动能守恒的。

在\Refx{公式}{完全弹性碰撞的计算公式}的右式两侧同乘以$(m_1+m_2)$，再平方，得到
\begin{Equation}[完全弹性碰撞1]
    \begin{cases}
        v_{1b}^2(m_1+m_2)^2=\qty[(m_1-m_2)v_{1a}+2m_2v_{2a}]^2=A\\
        v_{2b}^2(m_2-m_1)^2=\qty[(m_2-m_1)v_{2a}+2m_1v_{1a}]^2=B\\
    \end{cases}
\end{Equation}
关注到式\ref{式:完全弹性碰撞1}右式分别可以展开为
\begin{align*}
    A&=(m_1^2-2m_1m_2+m_2^2)v_{1a}^2+4(m_1-m_2)m_2v_{1a}v_{2a}+4m_2^2v_{2a}^2\\
    B&=(m_1^2-2m_1m_2+m_2^2)v_{2a}^2-4(m_1-m_2)m_1v_{1a}v_{2a}+4m_1^2v_{1a}^2
\end{align*}
现对式\ref{式:完全弹性碰撞1}的右式计算$Am_1+Bm_2$，关注到上方$A,B$展开式的第二项在相加时相消
\begin{align*}
    &Am_1+Bm_2\\
    =&m_1(m_1^2-2m_1m_2+m_2^2)v_{1a}^2+m_2(m_1^2-2m_1m_2+m_2^2)v_{2a}^2+4m_1m_2v_{2a}^2+4m_2m_1^2\\
    =&m_1(m_1^2+2m_1m_2+m_2^2)v_{1a}^2+m_2(m_1^2+2m_1m_2+m_2^2)v_{2a}^2
\end{align*}
这就得到
\begin{Equation}[完全弹性碰撞2]
    Am_1+Bm_2=m_1(m_1+m_2)^2v_{1a}^2+m_2(m_1+m_2)^2v_{2a}^2
\end{Equation}
现对式\ref{式:完全弹性碰撞1}的左式计算$Am_1+Bm_2$，这是十分容易的
\begin{Equation}[完全弹性碰撞3]
    Am_1+Bm_2=m_1(m_1+m_2)^2v_{1b}^2+m_2(m_1+m_2)^2v_{2b}^2
\end{Equation}
联立式\ref{式:完全弹性碰撞2}和式\ref{式:完全弹性碰撞3}，消去$(m_1+m_2)^2$并除以$2$，得到
\begin{Equation}[完全弹性碰撞4]
    \frac{1}{2}m_1v_{1a}^2+\frac{1}{2}m_2v_{2a}^2=\frac{1}{2}m_1v_{1b}^2+\frac{1}{2}m_2v_{2b}^2
\end{Equation}
这就证明了$E_{ka}=E_{kb}$，即\textbf{在完全弹性碰撞中，动量守恒和动能守恒同时成立。}

\subsection{完全非弹性碰撞}
当恢复系数$e=0$时，此时分离速度为零，称为\uwave{完全非弹性碰撞}。

在\Refx{公式}{一般对心碰撞的计算公式}中的左式，令$e=0$得
\begin{BoxEquation}[完全非弹性碰撞的计算公式]
    v_{1b}=v_{2b}=\frac{m_1v_{1a}+m_2v_{2a}}{m_1+m_2}
\end{BoxEquation}
完全非弹性碰撞的特征，即两物体在碰撞后不再分离，以同一速度运动。

完全非弹性碰撞中的动能并不守恒，且事实上是动能损失最大的碰撞方式。

碰撞过程中动能被转化为其他形式的能量，例如子弹射入木块，动能消耗于木块受子弹挤压的形变，例如人跳上运动的汽车，动能被人和汽车间的摩擦力转化为热量。

\subsubsection{两类碰撞的对比}
完全非弹性碰撞和完全弹性碰撞是两种极端情形，可以对比如下
\begin{Table}{两类碰撞的对比}{|lcc|l|}
    \hline
    碰撞类型&恢复系数&物理效果&典型示例\\
    \hline
    完全弹性碰撞&$e=1$&物体相撞前后的动能守恒&
    \makecell*[l]
    {
        1. 金属小球相互撞击。\\
        2. 金属小球撞击墙壁反弹。\\
        3. 牛顿摆
    }\\ \hline
    完全非弹性碰撞&$e=0$&物体相撞后以同一速度运动&
    \makecell*[l]
    {
        1. 黏土小球碰撞后黏在一起。\\
        2. 步枪子弹射入木块。\\
        3. 人跳上运动的汽车。
    }\\ \hline
\end{Table}
对于恢复系数$e\in(0,1)$的一般对心碰撞，我们可以称之为\uwave{非弹性碰撞}。

\section{角动量与角动量守恒}

\subsection{角动量与角冲量}
我们将质点的位矢$\vb*{r}$与其动量$\vb*{P}$的叉积，定义为该质点的\uwave{角动量}，记作$\vb*{L}$
\begin{BoxDefinition}[角动量]
    \vb*{L}=m\vb*{r}\times\vb*{v}=\vb*{r}\times\vb*{P}
\end{BoxDefinition}
角动量是与参照系有关的量，即便是对于两个相对静止的参照系，虽然动量相同，但是参照系原点的不同选取，会导致位矢发生变化，从而使得角动量产生差异。

角动量也称为\uwave{动量矩}，其方向垂直于位矢和动量构成的平面，其大小$L=Pr\sin\theta$的意义是位矢和动量构成的平行四边形的面积，这也是将角动量称为动量矩的原因。

我们将质点的位矢$\vb*{r}$与其所受力$\vb*{F}$的叉积，定义为该力的\uwave{力矩}，记作$\vb*{M}$
\begin{BoxDefinition}[力矩]
    \vb*{M}=m\vb*{r}\times\vb*{a}=\vb*{r}\times\vb*{F}
\end{BoxDefinition}
尽管力是参照系无关的，但是力矩的定义包含位矢量，因此力矩是参照系有关的量。

力矩的方向垂直于位矢和力构成的平面，其大小为$M=Fr\sin\theta$。

力矩在国际单位制中的单位是\si{N\cdot m=kg\cdot m^2\cdot s^{-2}}，量纲是\si{[L^2MT^{-2}]}。

力矩与功共享相同的量纲，但是两者的物理意义完全不同，因此力矩不能用焦耳作为单位。

现在我们来研究力矩和角动量间的关系，对角动量$\vb*{L}$求导
\begin{equation*}
    \Fdif{\vb*{L}}{t}=\Fdif{t}(\vb*{r}\times\vb*{P})=\vb*{r}\times\Fdif{\vb*{P}}{t}+\Fdif{\vb*{r}}{t}\times\vb*{P}
\end{equation*}
根据\Refx{公式}{牛顿第二定律的动量形式}和\Refx{定义}{瞬时速度}
\begin{equation*}
    \Fdif{\vb*{L}}{t}=\vb*{r}\times\vb*{F}+\vb*{v}\times\vb*{P}
\end{equation*}
根据\Refx{定义}{动量}，速度$\vb*{v}$与动量$\vb*{P}$的方向相同，从而叉积为零矢量
\begin{BoxPrinciple}[力矩第二定律]
    质点的角动量对时间的变化率等同于合力矩，
    
    且角动量的方向与合力矩相同。
    \begin{BoxEq}
        \Fdif{\vb*{L}}{t}=\vb*{r}\times\vb*{F}
    \end{BoxEq}
\end{BoxPrinciple}
\Refx{定律}{力矩第二定律}阐明了力矩的物理意义，即\textbf{力矩是动量矩的变化率}，这与先前在第二章中的\Refx{定律}{牛顿第二定律}中\textbf{力是动量的变化率}是十分相似的。

\Refx{定律}{力矩第二定律}将在稍后直接推出角动量定理，

现在我们研究两个相互作用的质点间的力矩关系，考虑以下示意图
\begin{Tikz}{两个相互作用的质点间的力矩关系}
    \lPoint{O}{0,0}
    \lShift{X}{O}{4,0}[above][x]
    \lShift{Y}{O}{0,4}[right][y]
    \draw[->] (O)--(X);
    \draw[->] (O)--(Y);
    \lShift{A}{O}{63:3.6}[above][m_1]
    \lShift{B}{O}{18:2.9}[below][m_2]
    \lRatio{Ac}{A}{O}{0.4}[left][\vb*{r}_1]
    \lRatio{Bc}{B}{O}{0.5}[above][\vb*{r}_2]
    \lRatio{A'}{A}{B}{0.35}[above right][\vb*{f}_1]
    \lRatio{B'}{B}{A}{0.35}[right][\vb*{f}_2]
    \draw[-latex,gray,dashed] (O)--(A);
    \draw[-latex,gray,dashed] (O)--(B);
    \draw[-latex] (A)--(A');
    \draw[-latex] (B)--(B');
    \draw[fill=white] (A) circle(0.04);
    \draw[fill=white] (B) circle(0.04);
\end{Tikz}
根据\Refx{定律}{牛顿第三定律}
\begin{equation*}
    \vb*{f}_1=-\vb*{f}_2
\end{equation*}
由于力$\vb*{f}_1$与$\vb*{r}_2-\vb*{r}_1$方向相同，同时力$\vb*{f}_2$与$\vb*{r}_1-\vb*{r}_2$方向相同，并考虑上式
\begin{equation*}
    \vb*{f}_1=\lambda(\vb*{r}_2-\vb*{r}_1)\qquad\vb*{f}_2=\lambda(\vb*{r}_1-\vb*{r}_2)
\end{equation*}
根据\Refx{定义}{力矩}，质点$\vb*{m}_1$受到的力矩$\vb*{M}_1$
\begin{equation*}
    \vb*{M}_1=\vb*{r}_1\times\vb*{f}_1=\lambda(\vb*{r}_1\times\vb*{r}_2-\vb*{r}_1\times\vb*{r}_1)=\lambda\vb*{r}_1\times\vb*{r}_2
\end{equation*}
根据\Refx{定义}{力矩}，质点$\vb*{m}_2$受到的力矩$\vb*{M}_2$
\begin{equation*}
    \vb*{M}_2=\vb*{r}_2\times\vb*{f}_2=\lambda(\vb*{r}_2\times\vb*{r}_1-\vb*{r}_2\times\vb*{r}_2)=\lambda\vb*{r}_2\times\vb*{r}_1
\end{equation*}
考虑到叉积的反交换律$\vb*{r}_1\times\vb*{r}_2=-\vb*{r}_2\times\vb*{r}_1$，由上两式得到
\begin{BoxPrinciple}[力矩第三定律]
    两个质点之间的相互作用力的力矩，大小相等，方向相反，分别作用在两个质点上。
    \begin{BoxEq}
        \vb*{M}_1=-\vb*{M}_2
    \end{BoxEq}
\end{BoxPrinciple}
\Refx{定律}{力矩第三定律}揭示了相互作用力的力矩关系，与\Refx{定律}{牛顿第三定律}是相似的。两者共同指出，质点之间的相互作用力与其力矩，均是大小相等方向相反的。

\Refx{定律}{力矩第三定律}将在稍后研究质点系的角动量定理中发挥重要作用。

如果有一个质点受到恒力矩$\vb*{M}$的作用，且这种作用持续了$\delt{t}$的时间，那么我们将力矩$\vb*{M}$与作用时间$\delt{t}$的乘积，称为力矩$\vb*{M}$产生的\uwave{角冲量}，记作$\vb*{H}$，这可以表示为
\begin{BoxDefinition}[恒力矩的角冲量]
    \vb*{H}=\vb*{M}\delt{t}
\end{BoxDefinition}

如果力矩$\vb*{F}$是随时间变化的变力矩，那么角冲量可以用定积分表示
\begin{BoxDefinition}[变力矩的角冲量]
    \vb*{H}=\Int[t_a][t_b]{\vb*{M}\dif{t}}
\end{BoxDefinition}
角冲量也称为\uwave{冲量矩}，是与参照系有关的量，因为力矩会随参照系而变。

根据\Refx{定义}{变力矩的角冲量}可知，\textbf{角冲量的物理意义是力矩在时间上的累积}。

我们知道，力矩是角动量的变化率，力矩对时间的积累即角动量的变化率对时间的积累，即在\Refx{定义}{变力矩的角冲量}中代入\Refx{定律}{力矩第二定律}
\begin{equation*}
    \vb*{H}=\Int[t_a][t_b]{\vb*{M}\dif{t}}=\Int[t_a][t_b]{\Fdif{\vb*{L}}{t}\dif{t}}=\Int[\vb*{L}_a][\vb*{L}_b]{\dif{\vb*{L}}}=\vb*{L}_b-\vb*{L}_a
\end{equation*}
这就得到了\uwave{角动量定理}
\begin{BoxPrinciple}[角动量定理]
    合力矩产生的角冲量，等于质点的角动量增量。
    \begin{BoxEq}
        \vb*{H}=\vb*{L}_b-\vb*{L}_a=\delt{\vb*{L}}
    \end{BoxEq}
\end{BoxPrinciple}
\Refx{定律}{角动量定理}是\Refx{定律}{力矩第二定律}的应用结果，因为力矩是角动量的变化率，因此角冲量，即力矩在时间上的积累，必然是角动量的变化量。

角动量和角冲量均不是绝对量，但是\textbf{角动量定理普遍适用于一切惯性系}。

角动量和角冲量在国际单位制中的的单位是\si{N\cdot m\cdot s=kg\cdot m^2\cdot s^{-1}}，量纲是\si{[L^2MT^{-1}]}。

\Refx{定律}{角动量定理}与先前的\Refx{定律}{动量定理}是相似的，两者均适用一切惯性系，两者分别指出，力矩在时间上的积累是角动量，力在时间上的积累是动量。

\Refx{定律}{角动量定理}指出，质点所受的合力矩为零时，质点的角动量守恒。

关于合力矩为零这一条件，存在两种可能性
\begin{enumerate}
    \item 所受的合力为零，此时的合力矩为零。
    \item 所受的合力与位矢的方向平行，使得叉积为零，导致合力矩为零。
\end{enumerate}
这就说明，合力为零仅是合力矩为零的充分条件，而不是必要条件，合力与位矢平行的情况是十分常见的，如果质点所受的力始终指向某一定点，那么，以该定点为原点的参照系中，该质点所受的力矩总是为零，该质点角动量守恒。这样指向定点的力称为有心力。
\begin{center}
    质点受有心力$\quad\Leftrightarrow\quad$质点角动量守恒$\quad\Leftrightarrow\quad$质点的合力矩为零
\end{center}

引力是典型的有心力，航天器绕地球飞行时，航天器仅受到地球的引力，如果以地球为原点建立参照系，那么航天器是角动量守恒的，取决于航天器的速度大小，航天器的运功轨迹可能是圆、椭圆、抛物线、双曲线。即质点角动量守恒时其运动形式是多样的，那么这些运动中是否存在某些一般规律？就像质点动量守恒时其总是静止或做匀速直线运动。

% 那么现在我们的问题就是，质点角动量守恒时，质点在运动学上是否有类似于\Refx{定律}{牛顿第一定律}的特征。

% 我们知道，动量守恒的运动学效果是质点静止或做匀速直线运动，角动量守恒的运动学效果则没有这么明显，下面我们来系统研究角动量守恒的运动学效果。

我们关注到，在引力作用下的运动轨迹均是平面曲线，这并不是引力的独特性质所致，而是角动量守恒的结果，根据\Refx{定义}{角动量}可知$\vb*{L}=m\vb*{r}\times\vb*{v}$，即$\vb*{L}$同时垂直于$\vb*{r},\vb*{v}$构成的平面，当$\vb*{L}$为常矢量时，意味着$\vb*{r},\vb*{v}$总在同一平面上，即质点是平面运动的。
\begin{BoxPrinciple}[质点角动量守恒的第一运动学效应]
    当质点的角动量守恒时，质点的运动轨迹均在同一平面上。
\end{BoxPrinciple}
\Refx{定律}{质点角动量守恒的第一运动学效应}是角动量守恒的必要条件，因此角动量守恒必然是平面运动，但是平面运动未必满足角动量守恒，例如对于一个在平面上的静止质点，现施加一个始终与其位矢垂直的力，此时力矩恒不为零，但是其运动始终在该平面上。

由于角动量守恒时$\vb*{L}$为常矢量，且方向是确定的，总是垂直于质点的运动平面，从而我们可以改用标量予以表示，即（下式的$\theta$是矢量$\vb*{r}$与矢量$\vb*{v}$间的夹角）
\begin{Equation}[角动量守恒第二效应1]
    L=mrv\sin\theta=mr\Fdif{r}{t}\sin\theta=m\frac{r\dif{r}\sin\theta}{\dif{t}}
\end{Equation}
考察一个$\dif{t}$时间内的微元过程，设$\dif{S}$为该过程中位矢$\dif{\vb*{r}}$扫过的面积
\begin{Tikz}{质点角动量守恒的第二运动学效应的示意图}
    \lPoint{O}{0,0}
    \lShift{X}{O}{4,0}[above][x]
    \lShift{Y}{O}{0,4}[left][y]
    \draw[->] (O)--(X);
    \draw[->] (O)--(Y);
    \lShift{A}{O}{73:3.6}
    \lShift{B}{O}{58:3.7}
    \lRatio{Ac}{A}{O}{0.4}[left][\vb*{r}]
    \lRatio{Bc}{B}{O}{0.5}[right][\vb*{r}+\dif{\vb*{r}}]
    \lRatio{V}{A}{B}{2.5}[below][\vb*{v}]
    \draw[-latex,gray,dashed] (O)--(A);
    \draw[-latex,gray,dashed] (O)--(B);
    \draw[-latex] (A)--(V);
    \draw[-latex] (A)--(B);
    \lRatio{C}{A}{B}{0.5}[above][\dif{\vb*{r}}]
    \draw[fill=white] (A) circle(0.04);
    \draw[fill=white] (B) circle(0.04);
    \fill[fill opacity=0.1,gray] (O)--(A)--(B)--cycle;
    \lRatio{S}{O}{B}{0.75}[left][\dif{S}]
    \lAngle[\theta]{O}{A}{B}
\end{Tikz}
由于$\vb*{v}$的方向与$\dif\vb*{r}$的方向一致，因此$\theta$也同时为$\vb*{r}$与$\dif\vb*{r}$的夹角。

根据正弦定理计算$\dif\vb*{S}$的面积
\begin{Equation}[角动量守恒第二效应2]
    \dif{S}=\frac{1}{2}r\dif{r}\sin\theta
\end{Equation}
上式代入式\ref{式:角动量守恒第二效应1}，得到
\begin{equation*}
    L=2m\Fdif{S}{t}
\end{equation*}
上式中的$\Fdif{S}{t}$可以视为位矢扫过的面积$S$随时间的变化率，即单位时间内位矢的扫掠面积，而角动量守恒使得$L$为一常量，这就可以得到一个重要的恒等式
\begin{BoxPrinciple}[质点角动量守恒的第二运动学效应]
    当质点的角动量守恒时，质点的位矢在单位时间内的扫掠面积为一常数。
    \begin{BoxEq}
        \Fdif{S}{t}=\frac{L}{2m}
    \end{BoxEq}
\end{BoxPrinciple}
这就得到了一个类似于\Refx{定律}{牛顿第一定律}的结论，一切物体在单位时间内的扫掠面积均为常数，直到有力矩迫使它改变这种状态为止，这就是力矩为零的运动学效应。

\Refx{定律}{质点角动量守恒的第二运动学效应}中，将质点视为行星，将质点所受的力视为行星受位于原点处的太阳引力，这就得到了开普勒第二定律，即行星在相同时间间隔内扫过的面积大小相同，即开普勒第二定律是该效应在有心力为引力时的特殊情形。

\subsection{质点系的角动量定理}
设两个质点的质量分别为$m_1,m_2$，在$t_a$至$t_b$的时刻间
\begin{enumerate}
    \item 两者分别受外力矩$\vb*{M}_{1,E},\vb*{M}_{2,E}$和内力矩$\vb*{M}_{1,I},\vb*{M}_{2,I}$的作用。
    \item 两者受到的外力矩的角冲量为$\vb*{H}_{1,E},\vb*{H}_{2,E}$，两者受到的内力矩的角冲量为$\vb*{H}_{1,I},\vb*{H}_{2,I}$。
\end{enumerate}
对质点$m_1$运用\Refx{定律}{角动量定理}，设$t_a,t_b$时其角动量分别为$\vb*{L}_{1,a},\vb*{L}_{1,b}$
\begin{Equation}[质点系角动量定理1]
    \vb*{H}_1=\vb*{H}_{1,E}+\vb*{H}_{1,I}=\Int[t_a][t_b]{(\vb*{M}_{1,E}+\vb*{M}_{1,I})\dif{t}}=\vb*{L}_{1,b}-\vb*{L}_{1,a}
\end{Equation}
对质点$m_2$运用\Refx{定律}{角动量定理}，设$t_a,t_b$时其角动量分别为$\vb*{L}_{2,a},\vb*{L}_{2,b}$
\begin{Equation}[质点系角动量定理2]
    \vb*{H}_2=\vb*{H}_{2,E}+\vb*{H}_{2,I}=\Int[t_a][t_b]{(\vb*{M}_{2,E}+\vb*{M}_{2,I})\dif{t}}=\vb*{L}_{2,b}-\vb*{L}_{2,a}
\end{Equation}
上两式成立，是因为质点所受合力矩的角冲量，等同于质点的内力矩和外力矩的角冲量和。
\begin{itemize}
    \item 外力矩对质点系的角冲量，等同于外力矩对各质点的冲量和，即$\vb*{H}_E=\vb*{H}_{1,E}+\vb*{H}_{2,E}$。
    \item 内力矩对质点系的角冲量，等同于内力矩对各质点的冲量和，即$\vb*{H}_I\hspace{2.2pt}=\vb*{H}_{1,I}\hspace{2.2pt}+\vb*{H}_{2,I}\hspace{2.2pt}$。
    \item 质点系的角动量，等同于各质点的角动量和，即$\vb*{L}_{a}=\vb*{L}_{1,a}+\vb*{L}_{2,a}$和$\vb*{L}_{b}=\vb*{L}_{1,b}+\vb*{L}_{2,b}$。
\end{itemize}
而我们知道，内力矩$\vb*{M}_{1,I},\vb*{M}_{2,I}$是一对相互作用力的力矩，根据\Refx{定律}{力矩第三定律}
\begin{equation*}
    \vb*{H}_{I}=\vb*{H}_{1,I}+\vb*{H}_{2,I}=\Int[t_a][t_b]{(\vb*{M}_{1,I}+\vb*{M}_{2,I})\dif{t}}=\vb*{0}
\end{equation*}
即质点系合内力矩的角冲量$\vb*{H}_I$恒为零矢量，因而式\ref{式:质点系角动量定理1}和式\ref{式:质点系角动量定理2}相加得到
\begin{align*}
    \vb*{H}_E
    &=\qty(\vb*{L}_{1,b}-\vb*{L}_{1,a})+\qty(\vb*{L}_{2,b}-\vb*{L}_{2,a})\\
    &=\qty(\vb*{L}_{1,b}+\vb*{L}_{2,b})-\qty(\vb*{L}_{1,a}+\vb*{L}_{1,b})\\
    &=\vb*{L}_{b}-\vb*{L}_{a}
\end{align*}\vspace{-20pt}
\begin{BoxPrinciple}[质点系的角动量定理]
    质点系的外力的角冲量的矢量和，等于质点系的角动量增量。
    \begin{BoxEq}
        \vb*{H}_E=\vb*{L}_{b}-\vb*{L}_{a}=\delt{\vb*{L}}
    \end{BoxEq}
\end{BoxPrinciple}
\Refx{定律}{质点系的角动量定理}的关键在于通过\Refx{定律}{力矩第三定律}论证了内力矩的角冲量和为零，从而说明了内力矩对系统的角动量无影响，这与在\Refx{定律}{质点系的动量定理}中应用\Refx{定律}{牛顿第三定律}以论证内力的动量和为零是类似的。

\Refx{定律}{质点系的角动量定理}只能在惯性系中使用，在非惯性系中不成立。

\subsection{角动量守恒定律}
\Refx{定律}{质点系的角动量定理}中，如果系统不受外力，关注到角动量不变。

这就得到了以下的\uwave{角动量守恒定律}
\begin{BoxPrinciple}[角动量守恒定律]
    当质点系是孤立的，所受的合外力矩为零，各质点仅受内力矩作用时，

    那么，各质点的角动量可以相互转化，但是系统的角动量总和保持不变，称为角动量守恒。
\end{BoxPrinciple}
角动量守恒定律指出，无论惯性系的原点在何处选取，对于孤立体系，其角动量总是守恒，我们知道对于地日系而言，如果以地球为研究对象，那么只有当惯性系的原点位于太阳时，地球才是角动量守恒的，相反，如果选取整个地日系为研究对象，那么无论惯性系的原点如何选取，尽管地球和太阳的角动量都在发生变化，但是整个地日系的角动量是守恒的。

\section{质心与质心运动定律}

\subsection{质心}
在讨论质点系的运动时，我们常引入质量中心的概念，称为\uwave{质心}。设质点系由$n$个质点组成，以$m_1,m_2,\cdots,m_n$分别表示各质点的质量，以$\vb*{r}_1,\vb*{r}_2,\cdots,\vb*{r}_n$分别表示各质点的位矢。

我们将质点系的总质量记为$m$，即
\begin{equation*}
    m=\SumIN{m_i}
\end{equation*}

我们将质心的位矢$\vb*{r}_C$定义为
\begin{BoxDefinition}[质心]
    \vb*{r}_C=\frac{1}{m}\SumIN{m_i\vb*{r}_i}
\end{BoxDefinition}
\Refx{定义}{质心}指出，质心的本质是质点位矢关于质量的加权平均。

\Refx{定义}{质心}适用于离散质点组成的质点系，如果质量分布是连续的，应改写为积分。
\begin{BoxGeneral}[连续体的质心][定义]
    如果刚体是一维细杆，线密度为$\lambda$
    \begin{BoxEq}
        \vb*{r}_C=\frac{1}{m}\Int[x_1][x_2]{\vb*{r}\lambda\dif{x}}
    \end{BoxEq}
    如果刚体是二维薄片，面密度为$\sigma$
    \begin{BoxEq}
        \vb*{r}_C=\frac{1}{m}\Idnt{\vb*{r}\sigma\dif{S}}
    \end{BoxEq}
    如果刚体是三维实体，体密度为$\rho$
    \begin{BoxEq}
        \vb*{r}_C=\frac{1}{m}\Itnt{\vb*{r}\rho\dif{V}}
    \end{BoxEq}
    以上的$r$均代表到原点的距离。
\end{BoxGeneral}
应当注意，质心未必位于连续体的内部，例如均匀圆筒的质心就位于其空心的中轴线处。

应当指出，质心与几何中心的概念是不同的，质心反映了质量的分布情况，只有当物体的密度是均匀分布时，质心与几何中心才是重合的，例如，均匀球体的质心位于球心处，但是如果球体一端密度较大一段密度较小，则球体的质心会偏向于密度较大的一端。

\subsection{质心运动定律}
当系统中每个质量都在运动时，系统质心的位置也会发生变化，现研究质心的动力学规律。

根据\Refx{定义}{质心}，质心的位矢是
\begin{equation*}
    \vb*{r}_C=\frac{1}{m}\SumIN{m_i\vb*{r}_i}
\end{equation*}
根据\Refx{定义}{瞬时速度}，质心的速度为
\begin{equation*}
    \vb*{v}_C=\Fdif{\vb*{r}_C}{t}=\frac{1}{m}\SumIN{m_i\Fdif{\vb*{r}_i}{t}}=\frac{1}{m}\SumIN{m_i\vb*{v}_i}
\end{equation*}
根据\Refx{定义}{瞬时加速度}，质心的速度为
\begin{equation*}
    \vb*{a}_C=\Fdif{\vb*{v}_C}{t}=\frac{1}{m}\SumIN{m_i\Fdif{\vb*{v}_i}{t}}=\frac{1}{m}\SumIN{m_i\vb*{a}_i}
\end{equation*}
将上式两侧同乘$m$，即得$m\vb*{a}_C=\SumIN{m_i\vb*{a}_i}$。

根据\Refx{定律}{牛顿第二定律}，设质点$m_i$所受的外力和内力分别为$\vb*{F}_{E,i},\vb*{F}_{I,i}$
\begin{equation*}
    m\vb*{a}_C=\SumIN{\vb*{F}_{E,i}+\vb*{F}_{I,i}}
\end{equation*}
根据\Refx{定律}{牛顿第三定律}，各质点内力和为零，记质点外力和为$\vb*{F}$
\begin{equation*}
    m\vb*{a}_C=\SumIN{\vb*{F}_{E,i}}=\vb*{F}
\end{equation*}
这就得到了\uwave{质心运动定律}
\begin{BoxPrinciple}[质心运动定律]
    无论物体的质量如何分布，无论外力作用在物体的什么位置，

    那么质心的运动，就等效于将物体的全部质量和全部外力集中于此处时质点的运动。
    \begin{BoxEq}
        \vb*{F}=m\vb*{a}_C
    \end{BoxEq}
\end{BoxPrinciple}
例如一枚手榴弹在抛物线轨道上爆炸时，尽管其碎片四处飞散，但是由于炮弹的爆炸力是内力，而内力不改变质心运动，因此事实上全部碎片的质心仍然依照抛物线运动。





\section{对称性与守恒律}
本章中我们得到了\Refx{定律}{动量守恒定律}和\Refx{定律}{角动量守恒定律}，在上一章末尾得到了\Refx{定律}{能量守恒定律}，尽管这三个结论都是在牛顿力学的框架下推证得出的，但在量子力学和相对论力学等更深层次的理论中，这些守恒定律仍然是成立的。

这就说明，这些守恒定律具有更普遍更深刻的物理基础，现代物理学的研究表明，三大守恒定律与自然界最为基本的属性——\textbf{时空对称性}是紧密联系在一起的。

这三大守恒律与时空对称性的联系可以阐述如下
\begin{BoxPrinciple}[守恒律与时空对称性]
    能量守恒，即时间平移对称性的体现，称为时间均匀性。\\[2mm]
    动量守恒，即空间平移对称性的体现，称为空间均匀性。\\[2mm]
    角动量守恒，即空间旋转对称性的体现，称为空间各向同性。\\[2mm]
\end{BoxPrinciple}
这就说明，相同的物理装置在相同条件下，无论装置放于何处，无论装置以何种朝向摆放，无论实验早些开始或晚些开始，物理实验的进程于结果，总是完全一样的。\footnote{这真是一个美好的世界，一切物理规律在时间和空间上竟然都是均匀的。}



